% Source: mse_optimal_predictor.tex
\begin{notechapter}
\chapter{The Optimal Predictor under Squared Loss}\label{ch:mse-optimal-predictor}

\section{Introduction}
In \textit{Deep Learning} by Goodfellow, Bengio and Courville (Section 6.2.1.2), the following result is derived using the calculus of variations:
\[
f^* = \arg\min_f \mathbb{E}_{x, y \sim p_{\text{data}}} \left[ \| y - f(x) \|^2 \right],
\]
where \( f^* \) is the optimal function that minimizes the expected mean squared error (MSE). The solution to this optimization problem is:
\[
f^*(x) = \mathbb{E}_{y \sim p_{\text{data}}(y \mid x)} [y].
\]
This document explains this result, its derivation, and its implications in detail.

\section{Key Concepts}

\subsection{Mean Squared Error (MSE)}
The MSE measures the average squared difference between the true output \( y \) and the predicted output \( f(x) \):
\[
\text{MSE} = \mathbb{E}_{x, y \sim p_{\text{data}}} \left[ \| y - f(x) \|^2 \right].
\]
Minimizing the MSE encourages the model to predict values close to the true values on average.

\subsection{Conditional Expectation}
The conditional expectation \( \mathbb{E}_{y \sim p_{\text{data}}(y \mid x)} [y] \) is the mean value of \( y \) given \( x \). It is the optimal predictor for \( y \) under the MSE criterion.

\section{Derivation of the Optimal Predictor}

\subsection{Decomposition of the MSE}
Using the law of total expectation, the expected MSE can be decomposed as:
\[
\mathbb{E}_{x, y \sim p_{\text{data}}} \left[ \| y - f(x) \|^2 \right] = \mathbb{E}_{x} \left[ \mathbb{E}_{y \sim p_{\text{data}}(y \mid x)} \left[ \| y - f(x) \|^2 \right] \right].
\]
This decomposition separates the overall MSE into an expectation over \( x \) of the conditional MSE given \( x \).

\subsection{Minimizing the Inner Expectation}
To minimize the overall MSE, we minimize the inner expectation for each \( x \):
\[
\mathbb{E}_{y \sim p_{\text{data}}(y \mid x)} \left[ \| y - f(x) \|^2 \right].
\]
This is the conditional MSE for a fixed \( x \). Minimizing it for each \( x \) ensures that the overall MSE is minimized.

\subsection{Calculus of Variations}
To find the minimizer, we take the derivative of the conditional MSE with respect to \( f(x) \) and set it to zero:
\[
\frac{d}{df(x)} \mathbb{E}_{y \sim p_{\text{data}}(y \mid x)} \left[ \| y - f(x) \|^2 \right] = 0.
\]
Expanding the derivative:
\[
-2 \mathbb{E}_{y \sim p_{\text{data}}(y \mid x)} [y - f(x)] = 0.
\]
Solving for \( f(x) \) gives:
\[
f(x) = \mathbb{E}_{y \sim p_{\text{data}}(y \mid x)} [y].
\]
Thus, the optimal function \( f^* \) is the conditional expectation of \( y \) given \( x \).

\subsection{A Derivation without Calculus}
The same result follows from a bias--variance decomposition. Writing
\( m(x) = \mathbb{E}[y \mid x] \) and expanding
\( y - f(x) = (y - m(x)) + (m(x) - f(x)) \):
\[
\mathbb{E}\left[ \| y - f(x) \|^2 \mid x \right]
= \underbrace{\mathbb{E}\left[ \| y - m(x) \|^2 \mid x \right]}_{\text{does not depend on } f}
+ \| m(x) - f(x) \|^2 ,
\]
because the cross term \( 2\,\mathbb{E}[y - m(x) \mid x]^\top (m(x) - f(x)) \) vanishes. The
right-hand side is minimized exactly when \( f(x) = m(x) \), and the minimal risk is the
expected conditional variance \( \mathbb{E}_x\big[\mathbb{E}[\| y - m(x) \|^2 \mid x]\big] \)
(the irreducible error).

\section{Uniqueness of the Solution}

\subsection{Uniqueness of the Conditional Expectation}
The conditional expectation \( \mathbb{E}_{y \sim p_{\text{data}}(y \mid x)} [y] \) is the unique minimizer of the conditional MSE because:
\begin{itemize}
    \item The squared error \( \| y - f(x) \|^2 \) is strictly convex in \( f(x) \), ensuring a unique global minimum.
    \item No other value of \( f(x) \) can achieve the same minimum expected squared error.
    \item Uniqueness of \( f^* \) as a function holds \emph{almost everywhere}: \( f^* \) can be modified on a set of inputs of probability zero (e.g.\ where \( p_{\text{data}}(x) = 0 \)) without changing the risk.
\end{itemize}

\subsection{Restricted Function Classes}
If the function class \( \mathcal{F} \) over which we optimize does not include the true conditional expectation, the decomposition above shows that minimizing the MSE over \( \mathcal{F} \) amounts to finding the element of \( \mathcal{F} \) closest to \( m \) in the \( L^2(p_{\text{data}}(x)) \) sense: the best predictor in the class is the projection of \( m \) onto \( \mathcal{F} \). For example:
\begin{itemize}
    \item In linear regression with a nonlinear true relationship, the least-squares solution is the best linear approximation of \( m \). It is still unique (as a function) when \( \mathbb{E}[x x^\top] \) is invertible, but its parameters may not be unique otherwise (collinear features).
    \item For non-convex classes (e.g.\ neural networks), several different functions may achieve the same minimal MSE, and the solution depends on the specific function class and the data distribution.
\end{itemize}

\section{Practical Implications}
\begin{itemize}
    \item In practice, we approximate \( f^* \) using models like neural networks or linear regression.
    \item If the model class is rich enough, it can approximate \( f^* \) arbitrarily well.
    \item Minimizing the MSE is a common objective in regression tasks because it aims to predict the conditional mean of \( y \) given \( x \).
\end{itemize}

\section{Conclusion}
The optimal predictor \( f^* \) that minimizes the expected MSE is the conditional expectation \( \mathbb{E}_{y \sim p_{\text{data}}(y \mid x)} [y] \). This solution is unique if the function class includes the true conditional expectation. In practice, the choice of model class and data distribution determines how well we can approximate \( f^* \).
\end{notechapter}
